\section{From Model to Interface}
\label{sec:implementation}

Section~\ref{sec:contract} defined the relations that must remain consistent.
This section follows the same boundary from the implementation side and focuses
on what changes for a visitor: a retained target highlight, an observable
handoff, capability-specific guidance, and either a grounded answer or a
recoverable rejection.  We continue the cheese-booth scenario in which the
Reception Agent hands a question about \texttt{refrigerated\_counter1} to the
Product Specialist Agent.  For readability, we call the canonical
FunctionalMLDS V2 instance the \emph{room contract}, its
trace-map/runtime-context projection the \emph{runtime map}, and the model hash
or contract fingerprint the \emph{content hash}; Appendix
Table~\ref{tab:runtime-terms} maps these shorter terms to the repository
terminology.

\subsection{Model relations that drive interface state}

A Meta-Language-defined Structure (MLDS) is a machine-readable artifact whose
vocabulary and structural rules are defined by an explicit meta-language
specification \cite{fischer2025machinereadable}.  Each room-specific
FunctionalMLDS document is such an instance.  Its role at interaction time is
to provide one authoritative chain from the selected spatial entity, through
the responsible agent and permitted capability, to the runtime action and the
evidence required before an answer may be shown.

FunctionalMLDS reuses selected UML and EAST-ADL use-case, requirements, and
verification concepts as its scenario foundation \cite{eastadl2013spec}; it
does not claim complete UML or EAST-ADL conformance.  The runtime-relevant
extension adds stable spatial entities, agent responsibilities and handoffs,
situated capability uses, concrete actions, and expected evidence.  Appendix
Table~\ref{tab:functionalmlds-relations} gives the compact element-and-relation
map needed to follow the migration and coverage claims in RQ1.

The distinction between a reusable capability and one situated use of it is
central.  ``Answer grounded questions'' is a capability; ``the Product
Specialist Agent answers about \texttt{refrigerated\_counter1} in this step''
is a capability use.  The latter connects a visible target and responsible
agent to one executable path, allowing the interface to retain the counter,
show a provider change, or report that the requested operation is unavailable.

A conventional dialogue manager or finite-state machine could reproduce the
nominal behavior of this single walkthrough.  The difference appears when
scene bindings, agent roles, handoffs, and backend actions evolve in separate
artifacts.  Consider a drift case in which Unity still selects the refrigerated
counter, the role configuration assigns it to the Product Specialist Agent,
but a stale backend route invokes a generic Reception Agent action.  A state
machine can prevent that failure only if it also imports and synchronizes the
same target, provider, action, and evidence relations.  FunctionalMLDS makes
that synchronization an explicit, typed, and versioned contract.  The
user-perceivable benefit is not additional interface chrome: the inconsistent
turn is blocked, the selection is retained, and guidance can identify the
failed relation instead of presenting a plausible response from the wrong
execution path.  We do not compare implementation effort against a state
machine, so our claim is consistency and fault localization across layers, not
that the formal representation is always the simplest authoring mechanism.

\subsection{Prototype cue rationale and empirical status}

The contract requires distinguishable target, responsibility, capability, and
outcome states, but it does not prescribe their visual form.  The prototype
uses a minimal mapping guided by four design constraints: spatial locality,
event salience, limited persistent clutter, and actionable recovery.  The
target is highlighted in the scene because the referent is spatial and the cue
can remain co-located with it while the visitor composes a question.  A
transient provider or handoff status appears only when responsibility changes,
making the transition visible without occupying permanent scene space.  Text
is reserved for blocked or unresolved states because recovery requires an
actionable instruction.  Exact identifiers, bindings, and assertion records
remain developer-facing.  This mapping operationalizes system-status visibility
and recoverability rather than exposing the complete contract
\cite{amershi2019guidelines}.

These choices were not selected through a prior comparative design study,
pilot study, or documented heuristic inspection.  They are fixed prototype
decisions and therefore a design hypothesis, not an optimized cue set.
Textual or spoken explanations, a persistent responsibility panel, and an
explicit clarification dialogue remain plausible alternatives.  The user
study is the first empirical assessment of the implemented mapping and examines
target clarity, handoff understanding, perceived responsibility, guidance, and
overall ease; it does not compare those alternatives or establish that the
chosen cues optimally calibrate trust.

Figure~\ref{fig:interface-cue-states} makes the mapping concrete.  It presents
the same target--provider--action--evidence boundary as
Figure~\ref{fig:admission-boundary}, but at the user-visible moments produced by
the implementation rather than as a second gate diagram.

\FloatBarrier
\begin{figure*}[!t]
  \centering
  \begin{tikzpicture}[
    font=\scriptsize,
    panel/.style={
      draw,
      rounded corners=2pt,
      minimum width=50mm,
      minimum height=54mm,
      inner sep=0pt
    },
    viewport/.style={
      draw,
      rounded corners=1.5pt,
      minimum width=44mm,
      minimum height=28mm,
      inner sep=0pt
    },
    object/.style={
      draw,
      double,
      double distance=0.7pt,
      rounded corners=1pt,
      align=center,
      text width=27mm,
      minimum height=8mm,
      inner sep=0.8mm
    },
    agent/.style={
      draw,
      circle,
      align=center,
      minimum size=9mm,
      inner sep=0.4mm,
      font=\tiny
    },
    status/.style={
      draw,
      rounded corners=1.5pt,
      align=center,
      text width=42mm,
      minimum height=11mm,
      inner sep=1mm
    },
    blocked/.style={
      status,
      dashed
    },
    handoff/.style={-{Latex[length=2mm]},semithick},
    denied/.style={-{Latex[length=2mm]},dashed}
  ]
    \node[panel] (selection) at (0,0) {};
    \node[font=\bfseries] at (0,2.25) {(a) Target retained};
    \node[viewport] at (0,0.55) {};
    \node[agent] at (-1.45,1.25) {Reception};
    \node[object] at (0,0.20)
      {Refrigerated\\cheese counter};
    \node[font=\tiny,align=center] at (1.45,1.25)
      {Product\\specialist};
    \node[status] at (0,-1.75)
      {\textbf{Selected target}\\
       Highlight remains while the question is composed};

    \node[panel] (handoffpanel) at (5.6,0) {};
    \node[font=\bfseries] at (5.6,2.25) {(b) Responsibility change};
    \node[viewport] at (5.6,0.55) {};
    \node[agent] (reception) at (4.65,1.25) {Reception};
    \node[agent] (specialist) at (6.55,1.25)
      {Product\\specialist};
    \draw[handoff] (reception) -- node[above,font=\tiny] {handoff} (specialist);
    \node[object] at (5.6,0.20)
      {Refrigerated\\cheese counter};
    \node[status] at (5.6,-1.75)
      {\textbf{Provider changed}\\
       Reception Agent $\rightarrow$ Product Specialist Agent};

    \node[panel] (blockedpanel) at (11.2,0) {};
    \node[font=\bfseries] at (11.2,2.25) {(c) Scoped rejection};
    \node[viewport] at (11.2,0.55) {};
    \node[agent] (lounge) at (10.25,1.25) {Lounge\\host};
    \node[agent] (specialist2) at (12.15,1.25)
      {Product\\specialist};
    \draw[denied] (lounge) -- node[above,font=\tiny] {not one hop} (specialist2);
    \node[font=\bfseries] at (11.2,1.25) {$\times$};
    \node[object] at (11.2,0.20)
      {Refrigerated\\cheese counter};
    \node[blocked] at (11.2,-1.75)
      {\textbf{Request blocked}\\
       No permitted one-hop route; selection remains available};
  \end{tikzpicture}
  \caption{Illustrative interface-state mockup for the running scenario.
  The panels show the target, handoff, and rejection states produced by the
  same contract boundary as Figure~\ref{fig:admission-boundary}.  This is an
  explanatory reconstruction rather than a pixel-accurate screenshot; the
  labels paraphrase runtime state and are not prescribed by the formal model.}
  \Description{Three schematic interface panels show a refrigerated cheese
  counter remaining highlighted while a question is composed, a handoff from
  the Reception Agent to the Product Specialist Agent, and a blocked request
  from the Lounge Host in which the selected counter remains available for
  correction.}
  \label{fig:interface-cue-states}
  \label{fig:functionalmlds-core}
\end{figure*}

\FloatBarrier

\subsection{When consistency is checked}

The four validation mechanisms are not four sequential operations paid on
every conversational turn.  Before a room opens, serialization and model
checks reject malformed fields, broken references, non-unique source
identifiers, and contradictory relations.  Before an interaction is enabled,
the executable-profile check requires one complete
target--provider--capability--action path.  Per turn, runtime checks bind the
request and response to the pinned room contract.  Appendix
Table~\ref{tab:enforcement-layers} lists the detailed invariants and failure
boundaries.

Room objects use stable \texttt{id} and \texttt{sourceId} values rather than
display labels or engine object names.  Agents have corresponding source
identifiers and typed references to their capabilities, responsibilities,
grounded objects, and handoff targets.  The room contract remains
authoritative, while the client and backend receive only the runtime subset
they need.  A visible label can consequently change without silently changing
the target or provider addressed by the interaction.

During authoring, structured room data and frozen semantic artifacts are
transformed into entities, agent roles, scenarios, capability uses, runtime
bindings, assertions, and validation cases.  Language models may propose
semantics, role labels, handoffs, and knowledge, but deterministic stages assign
runtime identifiers and routes.  Atomic publication and content hashes keep
the room contract and its client and backend projections together: a stale
client, mismatching backend, or partially regenerated room fails during setup
instead of entering an interaction with inconsistent responsibilities.

\subsection{3D-client selection and visible state}

The evaluated client is implemented in Unity, but the binding principle is
engine-independent: every selectable scene object must map to exactly one
stable model entity.  A camera ray selects a registered collider, the client
marks it with a spatial highlight, and that target remains active while the
visitor composes the request.  The states \texttt{none}, \texttt{resolved}, and
\texttt{ambiguous} are explicit.  Only a resolved state enables a deictic
request; the other states remain visible and lead to guidance rather than an
inferred target.

A scene-object component binds each relevant collider to an exact entity and
source identifier.  At startup, the registry rejects duplicate or conflicting
bindings and missing target-specific contexts.  The request carries the
content hash, exact target identifiers, bounded ray and candidate data, input
modality, active agent, and utterance.  The backend repeats the authoritative
identity check rather than relying on the client alone.

The main walkthrough and Figure~\ref{fig:interface-cue-states} remain on the
refrigerated counter.  Appendix Figure~\ref{fig:unity-binding-capture} is a
separate frozen implementation capture of \texttt{brand\_panel3}; it verifies
the same collider-to-entity mechanism and is not a second running example.

The client bridge uses the selected entity and server-resolved provider to
locate one target-specific scenario context.  It accepts only the scenario,
capability, binding, and action identifiers confirmed by the backend; zero or
several exact contexts leave the interaction unresolved.

\subsection{Authoritative routing, execution, and recovery}

At session setup, the backend loads the room contract and runtime map, builds
the executable context, and pins the content hash.  For each request, it rejects
malformed, stale, ambiguous, or conflicting spatial payloads and derives the
selected asset's group and zone from the room contract rather than from
client-provided responsibility fields.

Responsibility resolution applies the asset--group--zone precedence from
Section~\ref{sec:contract}.  The active agent may answer locally or hand off
once along a declared edge.  A resolved provider change is returned as explicit
state for the handoff cue; the absence of a permitted provider yields
relation-specific guidance.  The backend then selects one exact runtime action
from the trusted scenario, provider, and target, so a generic endpoint match is
insufficient.

Before generation, the backend snapshots mutable session state.  Failures that
can be decided from the request stop before the response model and leave the
conversation unchanged.  Violations visible only after generation, such as an
undeclared handoff or mismatching returned identifier, restore the snapshot and
suppress the output.  Model-call success alone is therefore insufficient for
presentation.

\paragraph{Responsiveness boundary.}
Only target/provider/action lookup, schema checks, evidence comparison, and
snapshot bookkeeping lie on the per-turn contract path; authoring and room
validation occur before the conversation.  The current evaluation nevertheless
does not measure end-to-end interaction time.  Its only timing evidence is an
in-memory structural workload of 6.15--10.06\,ms for the contract
representation; it excludes file parsing, Unity, networking, and model
inference.  These values are not conversational latency, and we make no claim
about perceived responsiveness.  Full turn time and users' sensitivity to
delay remain to be evaluated for live spatial interaction.

\subsection{Evidence, feedback, and identifier strictness}

The runtime trace serves two audiences.  For presentation, the client compares
the returned target, provider, route, binding, and action with its pending
selection and the admitted contract path.  Matching evidence permits the
answer to be shown while retaining the responsible-agent state; a mismatch
produces rollback or corrective guidance.  For diagnosis, runtime events also
retain the content hash, scenario and step, capability use, capability,
assertion references, and observation verdicts.  Missing observations remain
\texttt{inconclusive} rather than silently passing.

Strictness applies to structured contract metadata, not to the wording of the
generated answer.  The model may describe the counter in different natural
language, but the returned target, provider, binding, and action identifiers
must match the admitted path.  A known alias could be normalized before this
comparison; an unknown or different identifier is rejected even when the prose
appears plausible.  This deliberately favors referential and authorization
consistency over fuzzy identifier recovery.  The current prototype does not
implement approximate matching.

The visitor interface exposes the selected or unresolved target, makes a
provider change observable through the evaluated handoff state, and shows
concise status or guidance when the request cannot proceed.  Exact route
reasons, bindings, identifiers, and assertion records remain in developer logs.
The user study evaluates this lightweight mapping; richer explanations and
persistent responsibility panels are not implemented or compared.

\subsection{Scope and scaling boundary}

The contribution is not object anchoring or dialogue state by itself.
Reference-resolution and conversational-XR systems can maintain a local
candidate referent or dialogue context
\cite{lemaignan2012grounding,delatorre2024llmr,carcangiu2025tellxr}.  The room
contract starts from that selected candidate and additionally constrains who
may respond, whether a handoff is permitted, which concrete action may run, and
which identifiers must return before presentation.  Section~\ref{sec:related}
positions this boundary against those systems in more detail.

The evaluated corpus contains 93 routable assets and four to six agents per
room.  The runtime map avoids scanning the full authoring model for each turn,
and the interface exposes only the current target, provider transition, and
status rather than all modeled relations.  Thus the visible cue count does not
grow with the number of objects in the room.  We have not evaluated larger
agent graphs, multi-hop execution, crowded selections, or whether users remain
able to understand responsibilities in substantially denser scenes.

\subsection{End-to-end walkthrough}

The complete path is concrete.  The visitor selects the refrigerated cheese
counter and keeps its highlight while asking, ``What cheese is displayed
here?''  The client submits the counter's exact identifiers and the active
Reception Agent.  The backend resolves the Product Specialist Agent, verifies
the declared direct handoff, and invokes only the grounded-answer action
declared for that target and provider.  The response is shown only when its
evidence still names the same counter, provider, route, and action; otherwise
the provisional state is restored and guidance is displayed.  Starting the
same request from the Lounge Host requires two handoff edges through the
Reception Agent, so it is rejected before generation and the conversation
remains unchanged.
